\subsection{Definition of the renormalized spectral stress tensor}
  \label{subsec:spectral-stress}

  The renormalized projective entropy functional
  $S_\Pi^{\mathrm ren}[g;\mu]$ defines a covariant stress tensor by
  \begin{equation}
    \mathcal{T}^{\Pi,{\mathrm ren}}_{\mu\nu}
    \equiv
    -\frac{2}{\sqrt{g}}
    \frac{\delta S_\Pi^{\mathrm ren}}{\delta g^{\mu\nu}} .
  \end{equation}

  Formally, prior to renormalization, the first variation reads
  \begin{equation}
    \delta S_\Pi
    =
    \frac12 \Tr\!\left(A_g^{-1}\delta A_g\right),
  \end{equation}
  so that the response is governed by the resolvent kernel.
  After renormalization, local counterterms contribute additional
  geometric variations.

  In four dimensions, and for the minimal Laplacian $\xi=0$,
  the metric variation admits the general decomposition
  \begin{equation}
    \frac{\delta S_\Pi^{\mathrm ren}}{\delta g^{\mu\nu}}
    =
    c_{\mathrm EH}(\mu)\, G_{\mu\nu}
    +
    c_\Lambda(\mu)\, g_{\mu\nu}
    +
    \beta(\mu)\, B_{\mu\nu}
    +
    T^{\mathrm non\mbox{-}loc}_{\mu\nu}
    +
    A_{\mu\nu}.
    \label{eq:full-decomposition}
  \end{equation}

  Each term has a distinct geometric origin.
